% Options for packages loaded elsewhere
\PassOptionsToPackage{unicode}{hyperref}
\PassOptionsToPackage{hyphens}{url}
\documentclass[
  a4paper,
]{ltjsarticle}
\usepackage{xcolor}
\usepackage[margin=25mm]{geometry}
\usepackage{amsmath,amssymb}
\setcounter{secnumdepth}{-\maxdimen} % remove section numbering
\usepackage{iftex}
\ifPDFTeX
  \usepackage[T1]{fontenc}
  \usepackage[utf8]{inputenc}
  \usepackage{textcomp} % provide euro and other symbols
\else % if luatex or xetex
  \usepackage{unicode-math} % this also loads fontspec
  \defaultfontfeatures{Scale=MatchLowercase}
  \defaultfontfeatures[\rmfamily]{Ligatures=TeX,Scale=1}
\fi
\usepackage{lmodern}
\ifPDFTeX\else
  % xetex/luatex font selection
\fi
% Use upquote if available, for straight quotes in verbatim environments
\IfFileExists{upquote.sty}{\usepackage{upquote}}{}
\IfFileExists{microtype.sty}{% use microtype if available
  \usepackage[]{microtype}
  \UseMicrotypeSet[protrusion]{basicmath} % disable protrusion for tt fonts
}{}
\makeatletter
\@ifundefined{KOMAClassName}{% if non-KOMA class
  \IfFileExists{parskip.sty}{%
    \usepackage{parskip}
  }{% else
    \setlength{\parindent}{0pt}
    \setlength{\parskip}{6pt plus 2pt minus 1pt}}
}{% if KOMA class
  \KOMAoptions{parskip=half}}
\makeatother
\setlength{\emergencystretch}{3em} % prevent overfull lines
\providecommand{\tightlist}{%
  \setlength{\itemsep}{0pt}\setlength{\parskip}{0pt}}
\usepackage{bookmark}
\IfFileExists{xurl.sty}{\usepackage{xurl}}{} % add URL line breaks if available
\urlstyle{same}
\hypersetup{
  pdftitle={無名複素ゼロ閉包からの時空・波動計量の読出し ――最小読出し、無限閉包特殊解、対数時空写像、および写像後の \textbackslash lambda\textbackslash nu 不変性},
  pdfauthor={木原 範昭（Independent Researcher）　ORCID 0009-0004-6753-4020},
  hidelinks,
  pdfcreator={LaTeX via pandoc}}

\title{無名複素ゼロ閉包からの時空・波動計量の読出し
――最小読出し、無限閉包特殊解、対数時空写像、および写像後の
\(\lambda\nu\) 不変性}
\author{木原 範昭（Independent Researcher）　ORCID 0009-0004-6753-4020}
\date{2026年9月3日　v2.0　Version DOI 10.5281/zenodo.22282218}

\begin{document}
\maketitle

\subsection{要旨}\label{ux8981ux65e8}

本稿は、空間、時間、波長、振動数、速度をあらかじめ基礎変数として置かず、置換によって同一視された無名な複素数集合と複素二乗ゼロ閉包

\[
\sum_n z_n^2=0
\]

だけを出発点として、何が内部から読出し可能かを最小個数から順に調べる。

まず、複素数一個では二乗

\[
z^2=(a^2-b^2)+2iab
\]

という二次形式を形式的には持つが、共通複素スケールをゲージとして商した場合、一個だけからは非自明な相対量は得られないことを区別する。二個の非零複素数について

\[
z_1^2+z_2^2=0
\]

を課すと、

\[
z_2=\pm i z_1
\]

が従い、振幅比 \[1\] と相対位相 \[\pm\pi/2\]
が外部物差しなしに固定される。三個では、二乗した量が複素平面上の閉三角形を形成し、相対振幅比と相対位相の間に余弦則型の拘束が生じる。ここでは一般に連続自由度が残り、二乗ゼロ閉包だけから有限位相格子が自動的に生じるとは主張しない。

次に、本稿の中心的構成として、整数列 \[M_k\]
と各殻内の位相巡回からなる可算無限特殊解

\[
z_{k,n}
=
\frac{A}{M_k}
\exp\left[i\left(\phi_k+\frac{\pi n}{M_k}\right)\right],
\qquad
n=0,\ldots,M_k-1
\]

を与える。各殻は厳密に二乗ゼロ閉包し、さらに

\[
\sum_k \frac{1}{M_k}<\infty
\]

なら無限二乗和は絶対収束する。特に \[M_k=k(k+1)\] は、絶対収束と
\[M_{k+1}/M_k\to1\] を同時に満たす明示例である。

この特殊解では、無名性を保ったまま、殻の振幅から相対空間尺度

\[
L_k\propto M_k
\]

を、二乗位相集合の巡回次数から相対周期尺度

\[
P_k=M_k
\]

を読み出せる。したがって写像前には \[L_k/P_k\]
が一定であり、\[\nu_k^{(0)}:=1/P_k\] とすれば
\[L_k\nu_k^{(0)}=\mathrm{const}\] が成立する。

さらに、乗法的な尺度比を加法的な相対座標へ写す連続準同型は定数倍を除いて対数であることから、二つの殻
\[j,k\] の間に

\[
\Delta X_{jk}=C_X\ln\frac{L_j}{L_k},
\qquad
\Delta T_{jk}=C_T\ln\frac{P_j}{P_k}
\]

を定義する。この時空読出し後に、

\[
\lambda_{jk}:=|\Delta X_{jk}|,
\qquad
\tau_{jk}:=|\Delta T_{jk}|,
\qquad
\nu_{jk}:=\frac{1}{\tau_{jk}}
\]

と置けば、特殊解の \[L_k\propto P_k\] により

\[
\lambda_{jk}\nu_{jk}
=
\frac{C_X}{C_T}
=
\mathrm{const}
\]

が再び成立する。したがって本構成では、\[\lambda\nu\]
型の逆数関係は時空写像前だけでなく、実際に加法的な空間・時間尺度へ写像した後にも保存される。

最後に、同じ単位へ換算した二つの読出しを

\[
W=X+iT
\]

として一つの複素読出しにまとめると、元の二乗操作から

\[
dW^2=dX^2-dT^2+2i\,dX\,dT
\]

が得られ、その実部は \[1+1\] 次元 Lorentz
型二次形式となる。本稿は、この構成を現実の時空の一意な導出とは主張しない。示すのは、無名複素ゼロ閉包系に、相対角、相対尺度、周期尺度、加法的時空尺度、\[\lambda\nu\]
型不変量、および Lorentz
型符号を一貫した順序で読出せる明示的な無限特殊解が存在することである。

\textbf{キーワード}：無名複素数、二乗ゼロ閉包、読出し、創発時空、相対尺度、対数写像、波長、振動数、Lorentz
型計量、roots of unity

\begin{center}\rule{0.5\linewidth}{0.5pt}\end{center}

\subsection{1.
問題設定――値ではなく「何が読めるか」}\label{ux554fux984cux8a2dux5b9aux5024ux3067ux306fux306aux304fux4f55ux304cux8aadux3081ux308bux304b}

通常の波動論では、空間座標 \[x\]、時間 \[t\]、波長 \[\lambda\]、振動数
\[\nu\] を先に置き、その上で波を

\[
\exp\{i(kx-\omega t)\}
\]

のように記述する。本稿では順序を逆にする。

最初に与えるのは複素数の集合と、二乗和の閉塞だけである。

\[
\mathcal Z_N
=
\{z_1,\ldots,z_N\}/S_N,
\]

\[
Q(\mathcal Z_N)
=
\sum_{n=1}^{N}z_n^2
=0.
\tag{1}
\]

\(S_N\)
で商するのは、各複素数に物理的な名前または順序を与えないためである。計算上の添字は記憶位置にすぎず、読出しは置換に依存してはならない。

また式 (1) は、任意の \[\gamma\in\mathbb C^\times\] による共通変換

\[
z_n\mapsto \gamma z_n
\]

の下で

\[
Q\mapsto \gamma^2Q=0
\]

を保つ。したがって、式 (1)
だけでは共通振幅と共通位相は固定されない。本稿では、この共通自由度に依存する量を絶対読出しとは認めず、差、比、閉包次数、置換不変な対称性などを内部読出しの候補とする。

本稿の問いは単純である。

\begin{quote}
背景空間も背景時間もない無名複素数集合から、外部の物差しを輸入せず、どの段階で角、尺度、周期、空間、時間、波長、振動数に相当する相対量を構成できるか。
\end{quote}

\begin{center}\rule{0.5\linewidth}{0.5pt}\end{center}

\subsection{2.
一個の複素数――形式的二乗読出しとゲージ不変読出しの区別}\label{ux4e00ux500bux306eux8907ux7d20ux6570ux5f62ux5f0fux7684ux4e8cux4e57ux8aadux51faux3057ux3068ux30b2ux30fcux30b8ux4e0dux5909ux8aadux51faux3057ux306eux533aux5225}

一個の複素数を

\[
z=a+ib
\]

とする。二乗は

\[
z^2=(a^2-b^2)+2iab.
\tag{2}
\]

したがって、固定した複素代数基底 \[(1,i)\] に対しては、

\begin{itemize}
\tightlist
\item
  差の二次形式 \[a^2-b^2\]
\item
  交差項 \[2ab\]
\end{itemize}

が一個の複素数の内部に形式的に存在する。

しかし、共通複素スケール
\[z\mapsto\gamma z\]、特に共通位相回転までゲージとして商するなら、\[a\]
と \[b\] の個別値や \[ab\]
の数値は一般には不変ではない。したがって本稿では、

\begin{enumerate}
\def\labelenumi{\arabic{enumi}.}
\tightlist
\item
  \textbf{代数的に二乗形式が存在すること}
\item
  \textbf{ゲージを商した後にも相対量として読めること}
\end{enumerate}

を区別する。

一個だけでは比較対象がないため、完全な \[\mathbb C^\times\]
ゲージの下で非零点は一つの軌道に属する。したがって、本稿で採用する厳しい意味の相対読出しは二個から始まる。

この区別は後で重要になる。Lorentz 型二次形式は最終段階で再び式 (2)
と同じ二乗構造から現れるが、そのときの \[X,T\]
はすでに相対読出しとして構成された量である。

\begin{center}\rule{0.5\linewidth}{0.5pt}\end{center}

\subsection{3.
二個――外部角度ゲージなしの最初の相対角}\label{ux4e8cux500bux5916ux90e8ux89d2ux5ea6ux30b2ux30fcux30b8ux306aux3057ux306eux6700ux521dux306eux76f8ux5bfeux89d2}

非零の二個 \[z_1,z_2\] が

\[
z_1^2+z_2^2=0
\tag{3}
\]

を満たすとする。すると

\[
z_2^2=-z_1^2
\]

であり、

\[
\left(\frac{z_2}{z_1}\right)^2=-1.
\]

したがって

\[
\boxed{
\frac{z_2}{z_1}=\pm i
}
\tag{4}
\]

である。

式 (4) は直ちに

\[
\boxed{
\frac{|z_2|}{|z_1|}=1
}
\tag{5}
\]

と

\[
\boxed{
\arg z_2-\arg z_1
\equiv\frac{\pi}{2}
\pmod{\pi}
}
\tag{6}
\]

を与える。

ここで絶対位相は定まっていない。二個を同時に回しても式 (3)
は保たれる。しかし、\textbf{相対位相としての直角}は共通回転では消えない。

したがって二個の非自明ゼロ閉包は、外部の分度器を先に置かず、相対振幅
\[1\] と相対角 \[\pi/2\] を固定する最小構成である。

なお、ここで得られた \[\pi/2\]
を一般の多体系の全位相が四値に量子化されることと混同してはならない。式
(3) が固定するのは二個閉包固有の相対角である。

\begin{center}\rule{0.5\linewidth}{0.5pt}\end{center}

\subsection{4.
三個――相対振幅と相対位相の結合}\label{ux4e09ux500bux76f8ux5bfeux632fux5e45ux3068ux76f8ux5bfeux4f4dux76f8ux306eux7d50ux5408}

次に、非零の三個

\[
x^2+y^2+z^2=0
\tag{7}
\]

を考える。

二乗した量を

\[
X=x^2,
\qquad
Y=y^2,
\qquad
Z=z^2
\]

と書けば、

\[
X+Y+Z=0.
\tag{8}
\]

すなわち \[X,Y,Z\] は複素平面上で閉じた三角形を作る。\[Z=-(X+Y)\]
だから、

\[
|Z|^2
=
|X|^2+|Y|^2
+2|X||Y|\cos(\arg X-\arg Y).
\]

\(|X|=|x|^2\)、\[\arg X=2\arg x\] を用いると、

\[
\boxed{
\cos 2(\arg x-\arg y)
=
\frac{|z|^4-|x|^4-|y|^4}
{2|x|^2|y|^2}
}
\tag{9}
\]

を得る。

式 (9)
の右辺は相対振幅だけ、左辺は相対位相だけで構成される。したがって三個で初めて、二個では
\[1\]
に潰れていた振幅比が非自明な形自由度を持ち、その形と相対位相が同じ閉包条件によって結ばれる。

本稿ではこれを\textbf{振幅―位相の共役的読出し}と呼ぶ。ただし「共役的」は、互いに一方を固定すれば他方が拘束されるという意味で用い、正準共役変数を証明したという意味ではない。

また、式 (9) は一般に連続な形パラメータを許す。したがって

\[
\sum_{j=1}^{3}z_j^2=0
\]

だけから、相対位相が \[45^\circ\]、\[90^\circ\]
などの有限格子に自動的に制限されるとは言えない。有限位相巡回は、以下の無限特殊解では\textbf{解の構造として明示的に実現する}。

\begin{center}\rule{0.5\linewidth}{0.5pt}\end{center}

\subsection{5.
可算無限ゼロ閉包の構成的特殊解}\label{ux53efux7b97ux7121ux9650ux30bcux30edux9589ux5305ux306eux69cbux6210ux7684ux7279ux6b8aux89e3}

\subsubsection{5.1 一般殻構成}\label{ux4e00ux822cux6bbbux69cbux6210}

整数列

\[
M_k\in\mathbb N,
\qquad
M_k\ge2,
\qquad
k=1,2,\ldots
\]

を取る。\[A\in\mathbb C^\times\] を共通スケール、\[\phi_k\in\mathbb R\]
を殻ごとの共通位相として、各 \[k\] に \[M_k\] 個の複素数

\[
\boxed{
z_{k,n}
=
\frac{A}{M_k}
\exp\left[i\left(\phi_k+\frac{\pi n}{M_k}\right)\right]
}
\qquad
(n=0,\ldots,M_k-1)
\tag{10}
\]

を置く。

二乗すると

\[
z_{k,n}^2
=
\frac{A^2e^{2i\phi_k}}{M_k^2}
\exp\left(\frac{2\pi i n}{M_k}\right).
\tag{11}
\]

したがって、\[M_k\] 次単位根の和から

\[
\sum_{n=0}^{M_k-1}z_{k,n}^2
=
\frac{A^2e^{2i\phi_k}}{M_k^2}
\sum_{n=0}^{M_k-1}
\exp\left(\frac{2\pi i n}{M_k}\right)
=0.
\tag{12}
\]

各殻が厳密に二乗ゼロ閉包する。

\subsubsection{5.2
無限和の絶対収束}\label{ux7121ux9650ux548cux306eux7d76ux5bfeux53ceux675f}

二乗項の絶対値は

\[
|z_{k,n}^2|
=
\frac{|A|^2}{M_k^2}
\]

だから、全体系について

\[
\sum_{k=1}^{\infty}
\sum_{n=0}^{M_k-1}|z_{k,n}^2|
=
|A|^2
\sum_{k=1}^{\infty}\frac{1}{M_k}.
\tag{13}
\]

よって

\[
\boxed{
\sum_{k=1}^{\infty}\frac{1}{M_k}<\infty
}
\tag{14}
\]

なら二乗和は絶対収束する。したがって無名集合として項を並べ替えても和は変わらず、

\[
\boxed{
\sum_{k,n}z_{k,n}^2=0
}
\tag{15}
\]

が厳密に成立する。

\subsubsection{5.3
最小の明示例}\label{ux6700ux5c0fux306eux660eux793aux4f8b}

本稿では

\[
\boxed{
M_k=k(k+1)
}
\tag{16}
\]

を代表例とする。このとき

\[
\sum_{k=1}^{\infty}\frac{1}{M_k}
=
\sum_{k=1}^{\infty}
\frac{1}{k(k+1)}
=1
\tag{17}
\]

である。

さらに

\[
\frac{M_{k+1}}{M_k}
=
\frac{(k+1)(k+2)}{k(k+1)}
=
\frac{k+2}{k}
\longrightarrow1.
\tag{18}
\]

したがって、この列は

\begin{enumerate}
\def\labelenumi{\arabic{enumi}.}
\tightlist
\item
  無限二乗和の絶対収束
\item
  殻の巡回次数 \[M_k\to\infty\]
\item
  隣接尺度比 \[M_{k+1}/M_k\to1\]
\end{enumerate}

を同時に満たす。

\begin{center}\rule{0.5\linewidth}{0.5pt}\end{center}

\subsection{6.
無名集合からの殻の再同定}\label{ux7121ux540dux96c6ux5408ux304bux3089ux306eux6bbbux306eux518dux540cux5b9a}

式 (10) の各殻では、全要素の振幅が

\[
|z_{k,n}|
=
\frac{|A|}{M_k}
\tag{19}
\]

で一致する。\[M_k\]
を相異なるように選べば、無名集合を入力しても\textbf{同振幅の多重度}によって殻を再構成できる。

すなわち、ある振幅 \[r_k\] を持つ要素数を \[N(r_k)\]
とすれば、特殊解では

\[
N(r_k)=M_k,
\qquad
r_k=\frac{|A|}{M_k}.
\tag{20}
\]

したがって

\[
\boxed{
r_kN(r_k)=|A|
}
\tag{21}
\]

が全殻で共通となる。

\(|A|\) 自体は共通スケールなので読めなくても、殻 \[j,k\] の比

\[
\frac{r_j}{r_k}
=
\frac{M_k}{M_j}
\tag{22}
\]

は読める。

一方、二乗した各殻の位相集合は

\[
\left\{
\exp\left(2i\phi_k+\frac{2\pi i n}{M_k}\right)
\right\}_{n=0}^{M_k-1}
\]

であり、正 \[M_k\] 角形をなす。その最小巡回生成元は

\[
U_k
=
\exp\left(\frac{2\pi i}{M_k}\right),
\]

\[
U_k^{M_k}=1.
\tag{23}
\]

したがって \[M_k\]
は、振幅多重度からも位相巡回次数からも同一に再構成される。

この「二つの独立な読出しが同じ整数 \[M_k\]
を返す」ことが、本特殊解の中心構造である。

\begin{center}\rule{0.5\linewidth}{0.5pt}\end{center}

\subsection{\texorpdfstring{7. 写像前の空間尺度・周期尺度と
\(\lambda\nu\)
読出し}{7. 写像前の空間尺度・周期尺度と \textbackslash lambda\textbackslash nu 読出し}}\label{ux5199ux50cfux524dux306eux7a7aux9593ux5c3aux5ea6ux5468ux671fux5c3aux5ea6ux3068-lambdanu-ux8aadux51faux3057}

\subsubsection{7.1
相対空間尺度}\label{ux76f8ux5bfeux7a7aux9593ux5c3aux5ea6}

絶対振幅は読まない。しかし振幅の逆数比は

\[
\frac{|z_k|^{-1}}{|z_j|^{-1}}
=
\frac{M_k}{M_j}
\]

である。そこで振幅から得られる空間型相対尺度 \[L_k\]
を、共通定数を除いて

\[
\boxed{
L_k\propto M_k
}
\tag{24}
\]

と読む。

これは「\[M_k\] メートル」と絶対長さを宣言することではない。読めるのは

\[
\boxed{
\frac{L_j}{L_k}=\frac{M_j}{M_k}
}
\tag{25}
\]

という比だけである。

\subsubsection{7.2
相対周期尺度}\label{ux76f8ux5bfeux5468ux671fux5c3aux5ea6}

同じ殻の二乗位相は最小回転 \[2\pi/M_k\] を \[M_k\]
回繰り返すと元へ戻る。したがって、周期型読出しを

\[
\boxed{
P_k:=M_k
}
\tag{26}
\]

とする。

ここでも秒は導入していない。\[P_k\]
は「一巡に必要な内部位相ステップ数」である。

その逆数を振動数型読出しとして

\[
\boxed{
\nu_k^{(0)}:=\frac{1}{P_k}
=\frac{1}{M_k}
}
\tag{27}
\]

と定義する。

\subsubsection{7.3
写像前の逆数関係}\label{ux5199ux50cfux524dux306eux9006ux6570ux95a2ux4fc2}

式 (24) と式 (27) より、単位規格化を除いて

\[
\boxed{
L_k\nu_k^{(0)}
=\mathrm{const}
}
\tag{28}
\]

が成立する。

しかし式 (28) だけでは不十分である。\[L_k\] と \[P_k\]
はまだ乗法的な尺度ラベルであり、通常の加法的な空間座標・時間座標へ写像した後にも同じ関係が保たれるかを検算しなければならない。

これを次節で行う。

\begin{center}\rule{0.5\linewidth}{0.5pt}\end{center}

\subsection{8.
乗法的尺度から加法的時空への対数読出し}\label{ux4e57ux6cd5ux7684ux5c3aux5ea6ux304bux3089ux52a0ux6cd5ux7684ux6642ux7a7aux3078ux306eux5bfeux6570ux8aadux51faux3057}

\subsubsection{8.1 なぜ対数か}\label{ux306aux305cux5bfeux6570ux304b}

相対尺度比は乗法的である。三つの殻 \[i,j,k\] に対して

\[
\frac{L_k}{L_i}
=
\frac{L_k}{L_j}\frac{L_j}{L_i}.
\tag{29}
\]

これを加法的な相対座標差

\[
X_{ki}=X_{kj}+X_{ji}
\]

へ写す連続関数 \[f\] は

\[
f(ab)=f(a)+f(b)
\]

を満たさなければならない。正の実数上で連続性を仮定すれば、その一般形は

\[
\boxed{
f(r)=C\ln r
}
\tag{30}
\]

である。

したがって、対数は任意に選んだ便利な変換ではなく、\textbf{乗法的尺度比を加法的相対距離へ変える連続準同型として定数倍を除いて一意}である。

\subsubsection{8.2 空間読出し}\label{ux7a7aux9593ux8aadux51faux3057}

二殻 \[j,k\] の空間型相対間隔を

\[
\boxed{
\Delta X_{jk}
:=
C_X\ln\frac{L_j}{L_k}
}
\tag{31}
\]

と定義する。特殊解では \[L_k\propto M_k\] なので

\[
\Delta X_{jk}
=
C_X\ln\frac{M_j}{M_k}.
\tag{32}
\]

絶対原点は不要であり、\[X\mapsto X+X_0\] は読出しを変えない。この意味で
\[X\] は一方向の相対アフィン座標である。

\subsubsection{8.3
時間尺度読出し}\label{ux6642ux9593ux5c3aux5ea6ux8aadux51faux3057}

周期についても比は乗法的である。

\[
\frac{P_k}{P_i}
=
\frac{P_k}{P_j}\frac{P_j}{P_i}.
\]

したがって、周期尺度を加法的な時間尺度差へ写す同じ論理から

\[
\boxed{
\Delta T_{jk}
:=
C_T\ln\frac{P_j}{P_k}
}
\tag{33}
\]

を得る。

特殊解では \[P_k=M_k\] なので

\[
\Delta T_{jk}
=
C_T\ln\frac{M_j}{M_k}.
\tag{34}
\]

ここで注意すべき点がある。静的な無名集合だけから直接得られるのは、各殻の巡回対称性
\[C_{M_k}\] と周期尺度 \[P_k\] である。式 (33)
は\textbf{異なる時計の周期尺度を比較する相対時間尺度}であり、時間の矢または無限回巻き数をすでに導出したものではない。殻内部の位相状態は
\[S^1\] 型の時計位相を与えるが、完全な \[T\in\mathbb R\]
の時間順序には追加の動力学または unwrap 規則が必要である。

本稿ではこの限定を明示したうえで、空間尺度と時間尺度の相対計量の整合性を検討する。

\begin{center}\rule{0.5\linewidth}{0.5pt}\end{center}

\subsection{\texorpdfstring{9. 時空写像後の
\(\lambda\nu=\mathrm{const}\)}{9. 時空写像後の \textbackslash lambda\textbackslash nu=\textbackslash mathrm\{const\}}}\label{ux6642ux7a7aux5199ux50cfux5f8cux306e-lambdanumathrmconst}

ここが本稿の主要な検算である。

写像後の二殻間空間尺度を

\[
\boxed{
\lambda_{jk}
:=
|\Delta X_{jk}|
}
\tag{35}
\]

と読む。同様に、写像後の周期尺度を

\[
\boxed{
\tau_{jk}
:=
|\Delta T_{jk}|
}
\tag{36}
\]

とし、その逆数を

\[
\boxed{
\nu_{jk}
:=
\frac{1}{\tau_{jk}}
}
\tag{37}
\]

と定義する。

特殊解では \[L_j/L_k=P_j/P_k=M_j/M_k\] なので、

\[
\lambda_{jk}
=
C_X\left|\ln\frac{M_j}{M_k}\right|,
\]

\[
\tau_{jk}
=
C_T\left|\ln\frac{M_j}{M_k}\right|.
\]

よって

\[
\boxed{
\lambda_{jk}\nu_{jk}
=
\frac{C_X}{C_T}
=:\,c_*
}
\tag{38}
\]

を得る。

規格化単位で \[C_X=C_T\] と取れば

\[
\boxed{
\lambda_{jk}\nu_{jk}=1.
}
\tag{39}
\]

次元を持つ単位へ戻す場合、\[c_*=C_X/C_T\]
は空間尺度と時間尺度を換算する共通定数となる。

重要なのは、式 (39) を写像前の

\[
M_k\times\frac{1}{M_k}=1
\]

だけから結論していないことである。実際に

\[
M
\longmapsto
X=C_X\ln M,
\]

\[
M=P
\longmapsto
T=C_T\ln P
\]

という加法的な時空尺度へ写した\textbf{後で}、その時空上の間隔から再定義した
\[\lambda\] と \[\nu\] が式 (38) を満たすことを確認している。

したがって本特殊解では、\[\lambda\nu\]
型関係は時空読出しによって破壊されず、読出し写像と可換である。

\begin{center}\rule{0.5\linewidth}{0.5pt}\end{center}

\subsection{10. Lorentz
型二次形式}\label{lorentz-ux578bux4e8cux6b21ux5f62ux5f0f}

空間尺度と時間尺度を同じ単位へ換算するため

\[
\widetilde T
:=
\frac{C_X}{C_T}T
\]

と置く。以下では記号を簡単にするため再び \[T\] と書く。

相対読出しを一つの複素量

\[
\boxed{
W=X+iT
}
\tag{40}
\]

にまとめる。

ここで、本稿の出発点と同じ「複素二乗」を適用すると

\[
\boxed{
dW^2
=dX^2-dT^2+2i\,dX\,dT
}
\tag{41}
\]

となる。

したがって実部は

\[
\boxed{
\operatorname{Re}(dW^2)
=dX^2-dT^2
}
\tag{42}
\]

であり、\[1+1\] 次元 Lorentz 型符号 \[(+,-)\] を持つ。

虚部は

\[
\boxed{
\operatorname{Im}(dW^2)
=2\,dX\,dT
}
\tag{43}
\]

であり、一個の複素数の二乗

\[
(a+ib)^2=a^2-b^2+2iab
\]

に現れた交差項と同型である。

さらに特殊解では、任意の二殻 \[j,k\]
に対して同じ尺度比が空間側と時間側へ読まれるため、

\[
\frac{\Delta X_{jk}}{C_X}
=
\frac{\Delta T_{jk}}{C_T}
=
\ln\frac{M_j}{M_k}
\]

が成立する。ここで \[C_X,C_T\]
は乗法的尺度から加法的座標への連続準同型を定める定数であり、その符号は各読出し軸の向きに対応する。基礎構成は軸の正方向を固定しないので、同じ絶対単位
\[|C_X|=|C_T|\] に規格化すると

\[
\Delta X_{jk}
=
\sigma\,\Delta T_{jk},
\qquad
\sigma=\pm1
\]

と書ける。したがって、式 (42) 上では任意の殻対について

\[
\boxed{
ds_{jk}^2
=
\Delta X_{jk}^2-\Delta T_{jk}^2
=0
}
\tag{44}
\]

が成立する。すなわち、本特殊解では一部に null
型の枝が存在するのではなく、\textbf{全ての殻対区間が null である}。

これは §9 の \[\lambda_{jk}\nu_{jk}=c_*\]
と同じ構造の別表現である。本特殊解では

\[
L_k\propto P_k
\]

であることが、写像後の

\[
\lambda\nu=\mathrm{const}
\]

と、同一単位での

\[
ds^2=0
\]

を同時に与える。したがって「振幅から読む空間尺度」と「位相巡回から読む周期尺度」の自己双対性、\[\lambda\nu\]
型不変量、および null
読出しは、本特殊解における一つの恒等構造の三つの表現である。

ここで主張しているのは、現実の Lorentz
時空の全物理が導出されたということではない。式 (42)
は、本稿で構成した二つの相対尺度を複素二乗読出しへ再投入すると Lorentz
型二次形式が自然に現れる、という代数的結果である。また、全殻対が null
になることは本特殊解の性質であり、一般の二乗ゼロ閉包解まで null
に限定されることを意味しない。完全な Lorentz
対称性、因果順序、ブーストの物理的実現、\[3+1\]
次元化は別問題として残る。

\begin{center}\rule{0.5\linewidth}{0.5pt}\end{center}

\subsection{11.
漸近的連続極限}\label{ux6f38ux8fd1ux7684ux9023ux7d9aux6975ux9650}

代表例 \[M_k=k(k+1)\] について、空間・時間尺度の隣接差は

\[
\Delta X_k
=
C_X\ln\frac{M_{k+1}}{M_k}
=
C_X\ln\frac{k+2}{k}
\longrightarrow0,
\tag{45}
\]

\[
\Delta T_k
=
C_T\ln\frac{k+2}{k}
\longrightarrow0.
\tag{46}
\]

この細分化は一様ではなく、読出し座標に依存して指数的に細かくなる。基準を固定して

\[
X_k=C_X\ln M_k
=C_X\ln[k(k+1)]
\]

と取ると、大きな \[k\] で

\[
X_k\sim2C_X\ln k,
\qquad
k\sim\exp\left(\frac{X_k}{2C_X}\right).
\]

一方、式 (45) は

\[
\Delta X_k
=C_X\ln\left(1+\frac{2}{k}\right)
\sim\frac{2C_X}{k}
\]

なので、座標 \[X\] における局所分解能は漸近的に

\[
\boxed{
\Delta X(X)
\sim
2C_X
\exp\left(-\frac{X}{2C_X}\right)
}
\tag{46a}
\]

となる。したがって本構成の連続極限は、最初から一様な実数格子を仮定するものではなく、対数読出し座標の進行とともに離散刻みが指数的に細分される\textbf{位置依存の漸近連続化}である。

また、各殻の二乗位相の最小刻みは

\[
\Delta\theta_k
=
\frac{2\pi}{M_k}
=
\frac{2\pi}{k(k+1)}
\longrightarrow0.
\tag{47}
\]

したがって大きな \[k\] では、

\begin{enumerate}
\def\labelenumi{\arabic{enumi}.}
\tightlist
\item
  空間尺度軸の隣接間隔
\item
  時間尺度軸の隣接間隔
\item
  時計位相の刻み
\end{enumerate}

が同時に任意に細かくなる。

本稿ではこれを\textbf{漸近的局所連続極限}と呼ぶ。離散集合が有限領域で最初から実数直線そのものになるわけではないが、\[k\to\infty\]
で相対分解能がゼロへ向かうため、連続計量を近似する局所極限を持つ。

一般に、本構成で同じ性質を持つ整数列には少なくとも

\[
M_k\to\infty,
\]

\[
\sum_k\frac{1}{M_k}<\infty,
\]

\[
\frac{M_{k+1}}{M_k}\to1
\tag{48}
\]

を要求すればよい。例えば \[M_k\] が整数化された \[k^p\] 型で \[p>1\]
なら、この三条件を満たす広い特殊解族が得られる。

\begin{center}\rule{0.5\linewidth}{0.5pt}\end{center}

\subsection{12.
本構成の一般化}\label{ux672cux69cbux6210ux306eux4e00ux822cux5316}

式 (10) の本質は、特定の \[k(k+1)\] ではない。殻 \[k\] に

\[
M_k
\]

個の位相巡回点を置き、その共通振幅を

\[
|z_k|
\propto
\frac{1}{M_k}
\tag{49}
\]

と結ぶことである。

この条件により、同じ整数 \[M_k\] が

\begin{itemize}
\tightlist
\item
  振幅から得る逆尺度
\item
  位相集合から得る巡回次数
\end{itemize}

の双方に現れる。

したがって一般化された十分条件は、

\[
\boxed{
\text{inverse-amplitude scale}
\ \propto\
\text{phase-cycle order}
}
\tag{50}
\]

である。

式 (50)
が成立すれば、乗法的な相対空間尺度と相対周期尺度は同じ比を持つため、対数写像後にも

\[
\lambda\nu=c_*
\]

が保存される。

これは本稿の特殊解を超えた、同型構成に対する十分条件である。ただし、式
(50) が必要条件であること、あるいは無名ゼロ閉包一般解から必ず式 (50)
が現れることは本稿では証明しない。

\begin{center}\rule{0.5\linewidth}{0.5pt}\end{center}

\subsection{13.
先行研究との位置関係}\label{ux5148ux884cux7814ux7a76ux3068ux306eux4f4dux7f6eux95a2ux4fc2}

本稿の各数学的部品には既知の近縁研究がある。ただし、本稿が採る読出しの順序とは出発点が異なる。

Penrose の twistor 理論は、複素射影幾何と null 構造を用いて Minkowski
時空と共形幾何を記述する代表的研究である {[}1{]}。しかし、そこでは
Minkowski
時空またはその複素化・共形構造が理論的背景として既に位置付けられている。本稿は、時空座標を入力せず、無名複素ゼロ閉包の内部読出しから相対時空尺度を構成する点で問題設定が異なる。

Nottale の scale relativity
は、尺度変換そのものに相対論的構造を導入し、対数尺度を用いた Lorentz
型の scale transformation と de Broglie 関係の一般化を議論した
{[}2{]}。本稿との近さは、乗法的尺度を対数座標で扱う点にある。一方、本稿では物理的長さまたは時空分解能を最初から与えず、無名複素集合の振幅多重度と位相巡回次数から比較尺度そのものを先に再構成する。

Englman と Yahalom は、時間依存波動関数について、位相と log modulus
の間の厳密な相反関係を解析性から導出した
{[}3{]}。この研究は複素量の位相と対数振幅が独立ではないことを示す重要な先例である。ただし、同研究では時間
\[t\] が波動関数 \[\psi(t)\]
の独立変数として最初から与えられている。本稿では、周期尺度と空間尺度を背景時間なしの無名集合から読出すことを目的とする。

本稿執筆時点の文献調査では、

\[
\text{anonymous complex set}
\to
\sum z_n^2=0
\to
\text{minimal relative readout}
\to
\text{explicit infinite zero-closure solution}
\to
(L,P)
\to
(X,T)
\to
\lambda\nu=\mathrm{const}
\to
\operatorname{Re}(dW^2)=dX^2-dT^2
\]

という一連の構成を同じ出発点から実行した先行研究は確認できなかった。ただし、これは文献不存在の完全証明ではなく、\textbf{本稿の知る限り}という限定を伴う。

\begin{center}\rule{0.5\linewidth}{0.5pt}\end{center}

\subsection{14.
定理として示した部分と、物理解釈として残る部分}\label{ux5b9aux7406ux3068ux3057ux3066ux793aux3057ux305fux90e8ux5206ux3068ux7269ux7406ux89e3ux91c8ux3068ux3057ux3066ux6b8bux308bux90e8ux5206}

本稿の主張を過大にしないため、数学的に直接示した部分と、読出し解釈として採用した部分を分ける。

\subsubsection{14.1
直接示した部分}\label{ux76f4ux63a5ux793aux3057ux305fux90e8ux5206}

\begin{enumerate}
\def\labelenumi{\arabic{enumi}.}
\tightlist
\item
  二個の非自明閉包 \[z_1^2+z_2^2=0\] から \[z_2/z_1=\pm i\] が従う。
\item
  三個閉包から式 (9) の振幅―位相拘束が従う。
\item
  式 (10) の殻は各 \[k\] で厳密に二乗ゼロ閉包する。
\item
  条件 \[\sum1/M_k<\infty\] の下で無限二乗和は絶対収束する。
\item
  殻振幅の多重度と二乗位相巡回次数の双方から同じ \[M_k\] が読める。
\item
  乗法群 \[\mathbb R_{>0}\] から加法群 \[\mathbb R\]
  への連続準同型は定数倍の \[\ln\] である。
\item
  \(L_k\propto P_k\) なら、対数時空写像後にも式 (38) の
  \[\lambda\nu=\mathrm{const}\] が成立する。
\item
  \(W=X+iT\) の二乗の実部は \[dX^2-dT^2\] である。
\end{enumerate}

\subsubsection{14.2
読出し解釈として採用した部分}\label{ux8aadux51faux3057ux89e3ux91c8ux3068ux3057ux3066ux63a1ux7528ux3057ux305fux90e8ux5206}

\begin{enumerate}
\def\labelenumi{\arabic{enumi}.}
\tightlist
\item
  振幅逆数比を空間型尺度 \[L\] と呼ぶこと。
\item
  位相巡回次数を周期型尺度 \[P\] と呼ぶこと。
\item
  \(\ln L\) と \[\ln P\]
  をそれぞれ相対空間尺度・相対時間尺度と解釈すること。
\item
  写像後の空間尺度差を波長 \[\lambda\]、時間尺度差の逆数を振動数 \[\nu\]
  と読むこと。
\item
  \(\operatorname{Re}(dW^2)\) を物理的時空線素候補と読むこと。
\end{enumerate}

これらは互いに整合し、同じ特殊解上で閉じる。しかし一意性または現実物理との同一性は今後の検証対象である。

\begin{center}\rule{0.5\linewidth}{0.5pt}\end{center}

\subsection{15.
限界と次の検証課題}\label{ux9650ux754cux3068ux6b21ux306eux691cux8a3cux8ab2ux984c}

本稿で未解決の問題は明確である。

\subsubsection{15.1
一般解ではない}\label{ux4e00ux822cux89e3ux3067ux306fux306aux3044}

式 (10) は強い対称性を持つ構成的特殊解である。無名複素ゼロ閉包

\[
\sum_n z_n^2=0
\]

の一般解が必ず殻構造または式 (50) を持つとは示していない。

\subsubsection{15.2
殻ごとに可約である}\label{ux6bbbux3054ux3068ux306bux53efux7d04ux3067ux3042ux308b}

各殻が単独でゼロ閉包するため、全体系は局所閉包の直和として可約である。より強い次の課題は、有限部分系へ分解できない既約な無限閉包でも同じ読出しが成立するかである。

\subsubsection{15.3
時間の矢は未導出である}\label{ux6642ux9593ux306eux77e2ux306fux672aux5c0eux51faux3067ux3042ux308b}

静的集合から得るのは位相巡回と周期尺度である。時計位相を \[S^1\] から
\[\mathbb R\] へ unwrap
する物理過程、時間順序、因果方向は本稿では導出しない。

\subsubsection{\texorpdfstring{15.4 \(1+1\) から \(3+1\)
への拡張は未解決である}{15.4 1+1 から 3+1 への拡張は未解決である}}\label{ux304bux3089-31-ux3078ux306eux62e1ux5f35ux306fux672aux89e3ux6c7aux3067ux3042ux308b}

本稿で構成したのは一つの空間尺度方向と一つの時間尺度方向である。三つの独立空間方向が無名複素集合からどの拘束で選択されるかは別問題である。

\subsubsection{15.5
物理定数の絶対値は出していない}\label{ux7269ux7406ux5b9aux6570ux306eux7d76ux5bfeux5024ux306fux51faux3057ux3066ux3044ux306aux3044}

式 (38) の \[c_*\]
は二つの読出し単位の換算定数である。本稿はその数値を既知の光速 \[c\]
と同一視しない。現実の波動伝播へ接続するには、独立な観測または動力学的検算が必要である。

これらの限界は、本稿の構成的存在証明を損なうものではない。むしろ、次の反証可能な問いを明確にする。

\begin{quote}
より一般または既約な無名ゼロ閉包においても、振幅由来の相対尺度と位相由来の周期尺度が同じ内部パラメータへ因数分解され、時空写像後の
\[\lambda\nu\] 不変性まで保たれるか。
\end{quote}

\begin{center}\rule{0.5\linewidth}{0.5pt}\end{center}

\subsection{16. 結論}\label{ux7d50ux8ad6}

本稿では「時空とは何か」から始めず、「外部の物差しなしに何が読めるか」から出発した。

二個の複素ゼロ閉包は

\[
z_2=\pm iz_1
\]

を通じて、最初のゲージ不変な相対振幅と相対角を固定する。三個閉包では、二乗平面の閉三角形を通じて相対振幅と相対位相が同じ拘束の中で結ばれる。

さらに、

\[
z_{k,n}
=
\frac{A}{M_k}
\exp\left[i\left(\phi_k+\frac{\pi n}{M_k}\right)\right]
\]

という可算無限特殊解を構成すると、各殻で厳密なゼロ閉包が成立し、\[\sum1/M_k<\infty\]
の下で無限和も絶対収束する。この集合では、無名性を保ったまま

\[
\text{inverse amplitude}
\longleftrightarrow M_k
\longleftrightarrow
\text{phase-cycle order}
\]

という二重読出しが成立する。

その結果、写像前には空間尺度型 \[L_k\] と周期型 \[P_k\] が

\[
L_k\propto P_k
\]

を満たし、さらに両者を乗法尺度から加法尺度へ対数写像すると、

\[
\Delta X
=C_X\ln\frac{L_j}{L_k},
\qquad
\Delta T
=C_T\ln\frac{P_j}{P_k}
\]

が得られる。

最も重要な結果は、この\textbf{時空読出し後}に再定義した波長・振動数型量についても

\[
\boxed{
\lambda\nu=c_*
}
\]

が保存されることである。

さらに

\[
W=X+iT
\]

に複素二乗読出しを再適用すれば、

\[
\boxed{
\operatorname{Re}(dW^2)
=dX^2-dT^2
}
\]

という Lorentz 型二次形式が得られる。

したがって本稿は、無名複素ゼロ閉包系について、

\[
\boxed{
\text{最小相対読出し}
\to
\text{振幅―位相結合}
\to
\text{無限閉包特殊解}
\to
\lambda\nu\text{ 型読出し}
\to
(X,T)\text{ 読出し}
\to
\lambda\nu\text{ の再現}
\to
\text{Lorentz 型符号}
}
\]

という一つの閉じた構成経路を与えた。

現実の時空を一意に導出したという主張ではない。しかし、背景時空を置かない複素ゼロ閉包の中に、通常の波動・時空計量と同型の関係を最後まで自己整合的に読出せる無限特殊解が少なくとも一つ存在することを示した点が、本稿の中心的結果である。

\begin{center}\rule{0.5\linewidth}{0.5pt}\end{center}

\subsection{参考文献}\label{ux53c2ux8003ux6587ux732e}

{[}1{]} R. Penrose, ``Twistor Algebra,'' \emph{Journal of Mathematical
Physics}, \textbf{8}(2), 345--366 (1967). DOI:
\href{https://doi.org/10.1063/1.1705200}{10.1063/1.1705200}.

{[}2{]} L. Nottale, ``The Theory of Scale Relativity,''
\emph{International Journal of Modern Physics A}, \textbf{7}(20),
4899--4936 (1992). DOI:
\href{https://doi.org/10.1142/S0217751X92002222}{10.1142/S0217751X92002222}.

{[}3{]} R. Englman and A. Yahalom, ``Reciprocity between moduli and
phases in time-dependent wave functions,'' \emph{Physical Review A},
\textbf{60}, 1802--1810 (1999). DOI:
\href{https://doi.org/10.1103/PhysRevA.60.1802}{10.1103/PhysRevA.60.1802}.

\begin{center}\rule{0.5\linewidth}{0.5pt}\end{center}

\subsection{付録
A：対数読出しの一意性}\label{ux4ed8ux9332-aux5bfeux6570ux8aadux51faux3057ux306eux4e00ux610fux6027}

正の尺度比 \[r>0\] を加法座標へ移す写像 \[f\] が

\[
f(rs)=f(r)+f(s)
\tag{A1}
\]

を満たすとする。\[r=e^u\] と置き、

\[
g(u):=f(e^u)
\]

と定義すると、

\[
g(u+v)=g(u)+g(v).
\tag{A2}
\]

\(f\)、したがって \[g\] に連続性を要求すれば Cauchy
の加法方程式の連続解は

\[
g(u)=Cu
\]

だけである。ゆえに

\[
\boxed{
f(r)=C\ln r.
}
\tag{A3}
\]

したがって、本稿の \[X,T\]
に用いた対数は、尺度比の乗法を座標差の加法へ写すという要求の下で定数倍を除いて一意である。

\begin{center}\rule{0.5\linewidth}{0.5pt}\end{center}

\subsection{\texorpdfstring{付録 B：代表列 \(M_k=k(k+1)\)
の三条件}{付録 B：代表列 M\_k=k(k+1) の三条件}}\label{ux4ed8ux9332-bux4ee3ux8868ux5217-m_kkk1-ux306eux4e09ux6761ux4ef6}

代表列

\[
M_k=k(k+1)
\]

について、

\subsubsection{B.1 絶対収束}\label{b.1-ux7d76ux5bfeux53ceux675f}

\[
\sum_{k=1}^{\infty}\frac{1}{M_k}
=
\sum_{k=1}^{\infty}
\left(\frac1k-\frac1{k+1}\right)
=1.
\]

\subsubsection{B.2 無限分解能}\label{b.2-ux7121ux9650ux5206ux89e3ux80fd}

\[
M_k\to\infty.
\]

したがって殻内の位相刻み

\[
\frac{2\pi}{M_k}\to0.
\]

\subsubsection{B.3
相対尺度の局所連続化}\label{b.3-ux76f8ux5bfeux5c3aux5ea6ux306eux5c40ux6240ux9023ux7d9aux5316}

\[
\frac{M_{k+1}}{M_k}
=1+\frac{2}{k}
\to1,
\]

したがって

\[
\ln\frac{M_{k+1}}{M_k}
\to0.
\]

以上から、\[M_k=k(k+1)\]
は無限閉包の絶対収束と、空間尺度・時間尺度・位相尺度の漸近的細分を同時に満たす最も単純な明示例の一つである。

\end{document}
